\chapter{מבוא ללמידת מכונה}

למידת מכונה היא ענף של בינה מלאכותית העוסק בפיתוח אלגוריתמים המאפשרים למחשבים ללמוד מנתונים.
המושג \textenglish{Machine Learning} הוצג לראשונה על ידי \textenglish{Arthur Samuel} בשנת 1959,
והפך לאחד התחומים המרכזיים במדעי המחשב המודרניים.

\section{סוגי למידת מכונה}

ישנם שלושה סוגים עיקריים של למידת מכונה: למידה מונחית (\textenglish{Supervised Learning}),
למידה בלתי מונחית (\textenglish{Unsupervised Learning}), ולמידה מחיזוקים
(\textenglish{Reinforcement Learning}).

\subsection{למידה מונחית}

בלמידה מונחית, המודל מקבל קבוצת אימון הכוללת זוגות קלט-פלט. המטרה היא ללמוד פונקציה
\(\hat{f}: \mathcal{X} \to \mathcal{Y}\) המקרבת את הפונקציה האמיתית \(f\).
פונקציית ההפסד הנפוצה ביותר היא שגיאת הריבועים הממוצעת:

\begin{equation}
    \mathcal{L}(\theta) = \frac{1}{n} \sum_{i=1}^{n} \left( y_i - \hat{y}_i \right)^2
\end{equation}

כאשר \(\theta\) מייצג את פרמטרי המודל, \(y_i\) הוא הערך האמיתי ו-\(\hat{y}_i\) הוא הערך החזוי.

\subsection{למידה בלתי מונחית}

בלמידה בלתי מונחית, המודל מנסה לגלות מבנים נסתרים בנתונים ללא תוויות. אלגוריתמים כמו
\textenglish{K-Means} ו-\textenglish{Principal Component Analysis} (ראשי תיבות: \textenglish{PCA})
הם דוגמאות נפוצות.

\subsection{למידה מחיזוקים}

למידה מחיזוקים (\textenglish{Reinforcement Learning}) עוסקת בסוכן הלומד לפעול בסביבה
כדי למקסם פרס מצטבר. המסגרת הפורמלית של \textenglish{Markov Decision Process} (\textenglish{MDP})
מגדירה את הבעיה כעזרת רביעיית \((S, A, P, R)\).

\section{רשתות נוירונים}

רשתות נוירונים עמוקות (\textenglish{Deep Neural Networks}, קיצור: \textenglish{DNN})
הן מחלקה של מודלים המורכבת משכבות מרובות של נוירונים מלאכותיים. הפונקציה הבסיסית של
כל נוירון היא:

\begin{equation}
    z = \sigma\!\left( \mathbf{w}^\top \mathbf{x} + b \right)
\end{equation}

כאשר \(\sigma\) היא פונקציית ההפעלה, \(\mathbf{w}\) הוא וקטור המשקלות, \(\mathbf{x}\) הוא
וקטור הקלט, ו-\(b\) הוא הסטייה (\textenglish{bias}).

דוגמה ליישום פשוט של שכבה צפופה (\textenglish{Dense Layer}) ב-\textenglish{Python}:

\begin{LTR}
\begin{verbatim}
import numpy as np

def dense_layer(x, W, b, activation='relu'):
    z = np.dot(W, x) + b
    if activation == 'relu':
        return np.maximum(0, z)
    elif activation == 'sigmoid':
        return 1 / (1 + np.exp(-z))
    return z
\end{verbatim}
\end{LTR}

\section{סיכום}

פרק זה הציג את היסודות של למידת מכונה, כולל הגדרות בסיסיות, סוגי למידה, ורשתות נוירונים.
בפרקים הבאים נעמיק בארכיטקטורות מתקדמות כגון \textenglish{Transformer} ומנגנוני תשומת הלב.
